\runninghead{Lift Up Your Heads}
\properlane{lift-up-your-heads}{Lift Up Your Heads: An Aging World Ends, and Those Who Do Not Love It Rejoice}

The Gospel describes one set of events and two responses to it. When the sun,
the moon and the stars give their signs and the sea roars, men will be
\latin{arescéntes \dots\ præ timóre et exspectatióne}, ``withering away for
fear and expectation of what shall come upon the whole world'' (Lk 21:26).
In the same circumstances the disciples are told \latin{respícite, et leváte
cápita vestra}, ``look up and lift up your heads, because your redemption is
at hand'' (v.~28). The question is how the same events can wither some and
lift others. This interpretation takes the signs as real blows on a world
that is passing away and finds the difference between the two responses in
what each person loves. It follows St Gregory the Great, who preached a
homily on this Gospel (\work{Homiliae in Evangelia} I.1), and St Cyril of
Alexandria, who expounded it in his commentary on St Luke (sermon 139).

\subsection{A world grown old}

Gregory begins with the purpose of the prophecy. The Lord, ``desiring to find
us ready, announces what evils will follow upon the aging world, \latin{ut nos
ab ejus amore compescat}, to restrain us from loving it'' (I.1.1). The signs
are foretold so that fear may teach what peace did not: ``if we will not fear
God in quiet, we may at least fear his approaching judgment when worn down by
blows.''

He then counts what his hearers had already seen. Christ had said a little
earlier in the chapter that nation would rise against nation, with great
earthquakes and pestilences (Lk 21:10--11). ``Some of all these things we see
already done; others we dread as coming soon.'' Nation rising against nation
``we see more in our own times than we read in books''; earthquakes that
overthrow cities are reported ``from other parts of the world''; pestilences
are unceasing. The signs in the sun and moon and stars ``we do not yet see
openly,'' but before Italy was given over to the sword of the heathen ``we saw
fiery battle lines in the sky.'' His conclusion is an argument from what had
been fulfilled: \latin{sequentium rerum certitudo est praeteritarum
exhibitio}, ``the showing of things past is the certainty of things to
follow'' (I.1.1).

Later in the homily Gregory compares the world to a human body. In youth the
body is vigorous, the chest strong, the neck thick; in old age the stature is
bent, the neck withered, the breath short, and ``even when illness is absent,
its very health is often a sickness.'' So the world ``in its earlier years
flourished as in youth,'' and ``now it is weighed down by its very old age.''
He draws the practical conclusion at once: \latin{Nolite ergo \dots\ diligere
quem videtis diu stare non posse}, ``Do not love what you see cannot stand
for long'' (I.1.5), and he cites St John: ``Love not the world, nor the things
which are in the world'' (1~Jn 2:15).

\subsection{The head is the mind}

With that premise Gregory turns to the command. The verse before it promises
the Son of man ``with great power and majesty,'' and Gregory reads that verse
against the first coming: \latin{In potestate et majestate visuri sunt quem in
humilitate positum audire noluerunt}, ``they will see in power and majesty him
whom they would not hear when he was placed in humility'' (I.1.2). Those
words, he says, were spoken against the reprobate; now ``the words are turned
to the consolation of the elect.'' He reads the lifted head as the mind:
\latin{levate capita, id est exhilarate corda}, ``lift up your heads, that is,
gladden your hearts, \latin{quia dum finitur mundus, cui amici non estis, prope
fit redemptio quam quaesistis}, because while the world of which you are not
friends is ending, the redemption you have sought draws near'' (I.1.3).
Scripture often puts the head for the mind, he explains, since as the members
are ruled by the head, thoughts are ordered by the mind. Those who love God are
told to rejoice at the world's end, \latin{quia videlicet eum quem amant mox
inveniunt, dum transit is quem non amaverunt}, ``because they soon find the
one they love, while the one they did not love passes away'' (I.1.3).

His reasoning turns on St James: ``Whosoever therefore will be a friend of
this world, becometh an enemy of God'' (Jas 4:4). Whoever does not rejoice as
the end approaches ``testifies that he is its friend.'' Grief at the world's
ruin belongs to those ``who have planted the roots of the heart in love of
it,'' who do not seek the life to come. Then he asks the question that makes
the argument his own: \latin{Quid est vita mortalis, nisi via?} ``What is
mortal life but a road? And consider, my brothers, what it is to grow weary
in the labour of the road, and yet not to want the road to end'' (I.1.3).
Just before that question he tells his hearers that, knowing the joys of the
heavenly homeland, \latin{festinare ad ea quantocius debemus}, ``we ought to
hasten to them as fast as we can,'' and should wish to reach them ``by a
shorter road.''

The Epistle, read earlier, speaks of a night that is ending: ``The night is
passed, and the day is at hand. Let us therefore cast off the works of
darkness'' (Rom 13:12). St John Chrysostom, preaching on the verse before,
compares Paul's readers to runners, who rouse themselves the more when they
near the end of the course (\work{Homilies on Romans} 24, on v.~11). In this
Mass Paul's night and Gregory's road ask the same of the hearer, to hurry as
the end draws near; Gregory adds that only the friend of the world grieves to
see it come.

Between the Epistle and the Gospel the choir sings the Alleluia verse from
Psalm 84: \latin{Osténde nobis, Dómine, misericórdiam tuam: et salutáre tuum
da nobis}, ``Shew us, O Lord, thy mercy; and grant us thy salvation.'' Sung
just before the Gospel, it is the prayer of the people Gregory describes, who
have sought redemption and are then told that it is near.

Christ's parable of the fig tree, Gregory says, shows that the world is to be
trodden underfoot and despised. Christ says that as summer is known from the
fruit of the trees, the kingdom of God is known to be near from these things,
and Gregory concludes from the words:
\latin{fructus mundi ruina est. Ad hoc enim crescit, ut cadat}, ``the fruit of the world is
its ruin. For it grows in order to fall; it buds so as to consume in disaster
whatever it has budded.'' The summer is the kingdom, ``because then the
clouds of our mourning pass, and the days of life shine with the brightness
of the eternal sun'' (I.1.3).

On the closing saying Gregory observes that nothing in nature is more
durable than heaven and earth, and nothing passes as quickly as a word, which
exists only while it is being spoken. Yet Christ says that heaven and earth
will pass and his words will not: \latin{sine mutabilitate manentes sententias
exprimit meus sermo qui transit}, ``my speech, which passes, expresses
judgments that remain without change'' (I.1.4).

\subsection{The redemption is of the body and of creation}

Gregory explains why the disciple may rejoice; Cyril of Alexandria explains
what the disciple is to look up for. He first marks where Luke's discourse
turns. After foretelling Jerusalem ``girt about with armies,'' Christ
``transfers His words from this subject to the time of the consummation.''
The signs are cosmic: ``creation begins, so to speak, to be changed, and
brings unendurable terrors upon the inhabitants of earth,'' and ``the
unendurable fear of those things that are coming will suffice for the
destruction of many'' (sermon 139). Cyril reads the withering of the Gospel
literally, as men dying of fear.

Then Christ comes, ``not secretly nor obscurely, but as God and Lord, in glory
such as becomes Deity; and will transform all things for the better. For He
will renew creation, and refashion the nature of man to that which it was at
the beginning.'' Earlier in the same sermon Cyril contrasts this coming with
the first: ``His second advent from heaven will not happen secretly as did
His coming at first, but will be illustrious and terrible.'' The redemption
for which the disciples lift their heads is the resurrection: ``the dead
shall rise, and this earthly and infirm body shall put off corruption, and
shall clothe itself with incorruption by Christ's gift.'' Cyril draws the
consequence from St Peter, ``what sort of persons ought we to be,'' and from
Christ's command three verses after the lesson, ``Be you therefore always
watching'' (2~Pet 3:11; Lk 21:36).

Cyril's renewal of creation has a support in St Cyril of Jerusalem, who
teaches his candidates for baptism from the parallel verses of Matthew.
``Let us not sorrow, as if we alone died; the stars also shall die; but
perhaps rise again. And the Lord rolleth up the heavens, not that He may
destroy them, but that He may raise them up again more beautiful''
(\work{Catechetical Lectures} 15.3). St Bede says the same of the heaven that
passes away in v.~33: it is the air, which the flood once reached and the
fire of judgment will reach, and \latin{coelum et terra et transibit et
erit}, ``heaven and earth will both pass away and be,'' cleansed of their
present form and kept in their nature (\work{In Lucae evangelium expositio}
VI, on 21:33).

The two Fathers agree on what the signs are, real, cosmic and terrible, and
on what they mean for the disciple, which is joy. Gregory supplies the moral
ground of that joy in the object of a person's love. Cyril supplies its
content: a body that will not decay and a world remade.

The Postcommunion asks that the mercy received at the altar prepare us for
\latin{reparatiónis nostræ ventúra solémnia}, ``the coming solemnities of our
restoration,'' which Schuster takes to be Christmas (p.~323). Cyril has said
what the restoration will be when Christ comes again: creation renewed and
man's nature made again as it was at the beginning. The Postcommunion is
prayed after a Gospel of that second coming; the birth it prepares to keep
would then begin the restoration that the coming completes.

\subsection{Present terrors and the nearness of the end}

Gregory reads the disasters of his own day as signs. In its final sections the
homily recalls a storm that had struck three days before: ``a sudden whirlwind
uprooted ancient trees, destroyed houses, and overturned churches from their
foundations,'' and many who went to bed well died in the night (I.1.5). The
invisible judge, Gregory says, ``moved the breath of the slightest wind,
stirred up the storm of a single cloud, and undermined the earth''; what will
he do \latin{cum per semetipsum venerit}, ``when he comes in person''?
(I.1.6).

St Augustine raises the strongest objection to this kind of reasoning. A
bishop, Hesychius, had written to him that the afflictions of their own time
fulfilled the Gospel's words about men withering for fear. Augustine replies
that the portents and wars of his day were no greater than those which the
histories of the nations record from before Christ came (Letter 199.34--35),
and that the evils of the time fell alike on the good and the wicked, on
Christ's people and the devil's (199.37). Of three people who give opinions
about the time, the one who says the Lord will come sooner
\latin{optabilius loquitur, sed periculose fallitur}, ``speaks more
desirably, but is dangerously deceived'' (199.54).

Gregory's homily qualifies its own inference. He does read the calamities of
his time as the approach of the end, but he calls such terrors
\latin{sequentis irae praecones}, ``heralds of the wrath to follow,'' and adds
that they are as unlike the last tribulation \latin{quantum a potentia
judicis persona praeconis distat}, ``as far as the person of a herald is
from the power of a judge'' (I.1.6). He uses the storm to show on a small
scale what the coming judge can do, and to make his hearers fear him now. Then
he tells them what to do: ``correct your life, change your ways, overcome
temptations by resisting them, and punish by weeping the sins you have
committed.'' His last sentence gives the reason: \latin{Adventum namque
aeterni judicis tanto securiores quandoque videbitis, quanto nunc
districtionem illius timendo praevenitis}, ``you will see the coming of the
eternal judge the more securely, the more you now forestall his severity by
fearing it'' (I.1.6).

Augustine's objection tells against Gregory's sense that the end was close at
hand, and it stands. What Gregory draws from the signs does not depend on
that sense: they wean the heart from a world that is passing, and they teach
fear of the judge while there is time to amend.

The Collect asks Christ to come and rescue us \latin{ab imminéntibus
peccatórum nostrórum perículis}, ``from the threatening dangers of our sins.''
The danger it names lies in the one who prays, as does the danger Gregory has
in view when he tells his hearers to amend. Gregory puts the remedy in tears
and changed ways. At the offering the Secret asks for the cleansing from the
holy gifts themselves, that they make us, \latin{poténti virtúte mundátos},
``cleansed by their mighty power,'' come purer to their source.

\subsection{The lifted soul and the waiting that is not confounded}

The Introit sings \latin{Ad te levávi ánimam meam}, ``To thee have I lifted up
my soul,'' and the Gospel commands \latin{leváte cápita vestra}. Gregory says
what the lifting of the head is: \latin{Levare itaque capita est mentes nostras
ad gaudia patriae coelestis erigere}, ``to lift up the head is to raise our
minds to the joys of the heavenly homeland'' (I.1.3). If that is the command,
the psalmist in the Introit's first line is already doing what the Gospel
asks. He lifts his soul while enemies threaten to laugh at him, as the
disciples lift their heads while the world is shaken.

The Gradual and the Offertory sing that \latin{univérsi, qui te exspéctant,
non confundéntur}, ``none who wait for thee shall be confounded.'' The Gospel
uses the noun of the same verb for the men who wither \latin{præ timóre et
exspectatióne}. The two kinds of waiting differ in what is waited for. The men
in the Gospel expect what will come upon the world, and Gregory has said that
grief at its ruin belongs to those who love it; those who wait for God wait
for him, and the psalm says they are never put to shame.

The Communion promises another harvest. Gregory says that the world's fruit is
its ruin; the psalm verse the Communion sings says that ``our earth shall yield
her fruit,'' and gives the reason first: \latin{Dóminus dabit benignitátem},
``The Lord will give goodness.''

\begin{fourSenses}
\item[Literal.] Christ foretells signs in the heavens and distress on earth,
his own coming in a cloud with power and majesty, and commands his disciples
to lift their heads because their redemption is near; heaven and earth will
pass, and his words will not (Lk 21:25--33).
\item[Allegorical.] Christ, who came first in secret and in humility and was
not heard, comes again openly, as God and Lord, in power and majesty (Cyril of
Alexandria; Gregory I.1.2); those who love him find him as the world they did
not love passes (Gregory I.1.3).
\item[Moral.] Do not love what cannot stand for long; a traveller who is
tired of the road should want it to end (Gregory). Fear the judge now, and
amend your life, so as to see his coming securely (Gregory). The danger to
pray against is that of one's own sins (Collect).
\item[Anagogical.] When the clouds of mourning pass, the days of life shine
with the eternal sun (Gregory); the dead rise and the body puts on
incorruption, and the heavens are rolled up to be raised more beautiful
(Cyril of Alexandria; Cyril of Jerusalem).
\end{fourSenses}
